\section{Full coloring in the random planted model} \label{app: plantedcoloring}

This part is dedicated to proving \cref{thm: random-planting-full}. We start with the key first step that guarantees the recovery of a \emph{list} of candidate $k$-colorings, such that at least one of them is close to the planted $k$-coloring.


\begin{theorem}[Result for random planting, partial recovery of a list]
    \label{thm: random-planting-partial-list}
    Let $H$ be an $n$-vertex $d$-regular graph with adjacency matrix $A_H$.
    Suppose $\rank_{\geq \lambda}(\frac{1}{d}A_H) \leq t$ for some $0 < \lambda < 1$ and $t \in \mathbb{N}$.
    Let $G$ be the graph obtained by randomly planting a \kcl in $H$, and let $\chi: [n] \to [k]$ be the associated \kcl.
    Then, for $\lambda > 0$ small enough, there exists an algorithm that runs in time $n^{O(1)} \cdot \Paren{O_{k}(1)}^{t}$ and outputs a list of $\Paren{O_{k}(1)}^{t}$ $k$-colorings $\hat{\chi}: [n] \to [k]$ such that with high probability, for at least one of them, it holds up to a permutation of the color classes that $\mathbb{E}_{\bm{x} \sim \operatorname{Unif}([n])} \1[\hat\chi(\bm{x}) \neq \chi(\bm{x})] \leq O_k(1/d^{\Omega(1)})$.
\end{theorem}
\begin{proof}
    The proof is identical to the proof of \cref{thm: random-planting-partial}, but omitting the last rounding step.
\end{proof}


To convert a coloring returned by \cref{thm: random-planting-partial-list} into a full coloring, we will need several helper lemmata. We start with the proofs of the two lemmata showing properties of \textbf{one}-sided expansion.

\begin{proof}[Proof of \cref{lemma: emlonesided}]
    Let $\1_S \in \{0,1\}^n$ be the indicator vector of $S$ and write $\1_S$ in the orthonormal eigenbasis $e_{[n]}$ of $A$:
    $$\1_S = \sum_{i=1}^n \alpha_i e_i\,.$$

    Hence, by orthonormality, we can write

    \begin{align*}
        &\Abs{E(H_S)}\\
        &=\frac{1}{2} \1_S^\top A \1_S\\
        &=\frac{1}{2} \Paren{\sum_{i=1}^n \alpha_i e_i^\top } \Paren{\sum_{i=1}^n \lambda_i(A) e_i e_i^\top } \Paren{\sum_{i=1}^n \alpha_i e_i}\\
        &=\frac{1}{2} \sum_{i=1}^n \lambda_i(A) \alpha_i^2
    \end{align*}

    Using $d$-regularity, we have $\lambda_1(A)=d$.
    Then, as $e_1=\frac{1}{\sqrt{n}}\1$, we get $\alpha_1=\Iprod{e_1, \1_S }=\frac{|S|}{\sqrt{n}}$, so $\lambda_1(A) \alpha_1^2=\frac{d|S|^2}{n}$.
    For the terms with $i>1$, as $\alpha_i^2\geq 0$, we can upper bound $\lambda_i(A) \alpha_i^2$ by $\lambda_2(A) \alpha_i^2\leq cd\alpha_i^2$.
    Plugging these back we get
    $$\Abs{E(H_S)}\leq \frac{1}{2}\Paren{\frac{d|S|^2}{n} + cd\sum_{i=2}^n \alpha_i^2}\,.$$

    Then, using $\sum_{i=2}^n \alpha_i^2\leq \sum_{i=1}^n \alpha_i^2=\Norm{\1_S}_2^2=|S|$ gives the desired bound.
\end{proof}

\begin{proof}[Proof of \cref{lemma: ssveonesided}]
    Let $T:=N(S)\cup S$.
    Then, as all edges touching $S$ are in $E(H_{T})$, we have
    $$e(H_{T})\geq \frac{1}{2}|S|d\,.$$ 

    On the other hand, by \cref{lemma: emlonesided}, writing $|T|=x|S|$ and using $|S|\leq \alpha n$, we get
    $$e(H_{T})\leq \frac{d|T|}{2}\Paren{\frac{|T|}{n}+c}\leq\frac{1}{2}x\Paren{c+\alpha x}|S|d\,.$$

    Comparing the two bounds, and solving the resulting quadratic inequality, we get 
    $$x\geq \frac{\sqrt{c^2 + 4\alpha}-c}{2\alpha}=\frac{2}{c+\sqrt{c^2+4\alpha}}\,.$$

    By using $\sqrt{a+b}\leq \sqrt{a}+\sqrt{b}$ and that $\Abs{N(S)\setminus S}\geq \Abs{N(S)}-|S|$ we get the desired result.
\end{proof}


Next, we need an algorithm that takes a coloring that agrees with most of the planted coloring and turns it into a coloring that leaves some of the vertices uncolored, but agrees with the planted coloring on the colored vertices. We define this notion of coloring formally now.

\begin{definition}[Partial coloring]
    \label{def: partialcoloring}
    Let $G$ be an $n$-vertex graph that admits a \kcl $\chi$.
    A {\em $(k, \beta)$-partial} coloring w.r.t.\ $\chi$ is a \kcl of at least $(1-\beta)n$ vertices that is identical to $\chi$ on these vertices.
    The remaining $\beta n$ vertices are left uncolored and are referred to as {\em free}. In the random planted model we always implicitly take $\chi$ to be the planted \kcl.
\end{definition}

Let us state now the corresponding algorithm and its guarantees.

\begin{algorithm}[H]
    As long as there is a vertex $x$ that has color $c_1$, and a color $c_2\neq c_1$ such that $x$ has less than $\frac{1}{6k}d$ neighbors with color $c_2$, uncolor $x$.%\andor{why does \cite{MR3536556-David16} need another condition?}
    \protect\caption{\label{alg: uncoloring}Uncoloring Algorithm}
\end{algorithm}


\begin{lemma}[Adaptation of Lemma B.20 of \cite{MR3536556-David16}]\label[lemma]{lemma: approximatetopartial}
    Let $\alpha\leq\frac{1}{24k}$ be a parameter. Let $H$ be a $d\geq O(1)$-regular $n$-vertex  graph with $\lambda_2(H)\leq \frac{1}{9k}$. 
    Let $G$ be the graph obtained from $H$ by randomly planting a \kcl $\chi$ in it. Given an \kcl $\hat{\chi}$ of $G$ with $\mathbb{E}_{\bm{x} \sim \operatorname{Unif}([n])} \1[\hat{\chi}(\bm{x}) \neq \chi(\bm{x})] \leq \alpha$, \cref{alg: uncoloring} \whp returns an $(k, 4\alpha)$-partial coloring of $G$.
\end{lemma}

To prove \cref{lemma: approximatetopartial}, we will need a definition and two lemmata: one upper bounding the number of atypical vertices with high probability, and a structural lemma showing the existence of a ``nice'' subgraph.


\begin{definition}\label{def: statisticallybad}
    Let $H$ be a $d$-regular graph. Let $G$ be the graph obtained from $H$ by randomly planting a \kcl $\chi$ in it. We call a vertex $x\in V(G)$ \textit{statistically bad}, if there is a color $c\neq \chi(x)$ such that $x$ has less than $\frac{1}{2k}d$ neighbors $y$ with $\chi(y)=c$. We denote the set of statistically bad vertices by $SB(G)$, writing $SB$ when it is clear from context which graph we are referring to.
\end{definition}

\begin{lemma}[Restatement of Lemma B.14 of \cite{MR3536556-David16}]\label[lemma]{lemma: sbupperbound}
   Let $H$ be a $d\geq O(1)$-regular $n$-vertex graph. Let $G$ be the graph obtained from $H$ by randomly planting a \kcl in it. It holds that $\Abs{SB(G)}\leq n2^{-\Omega(d)}$ \whp.
\end{lemma}

\begin{lemma}[Adaptation of Lemma B.19 of \cite{MR3536556-David16}]\label[lemma]{lemma: highmindegreesubgraph}
    Let $0<\alpha, c<1$ be parameters. Let $H$ be a $d$-regular $n$-vertex graph with $\lambda_2(H)\leq c$. For any $S\subseteq V$ with $\Abs{S}<\alpha n$, there exists $R\subseteq V\setminus S$ such that $\Abs{R}\geq \Paren{1-2\alpha}n$ and the minimum degree in $H_R$ is at least $(1-2\alpha-c)d$.
\end{lemma}

\begin{proof}
    Let $\e:=2\alpha + c$. Consider the following iterative procedure. Repeatedly remove a vertex $v_t$ from $V\setminus \Paren{S \cup \bigcup_{i=1}^{t-1} v_i}$ if $v_t$ has more than $\e d$ neighbors in $S_t := S \cup \bigcup_{i=1}^{t-1} v_i$. Assume towards a contradiction that this process stops after more than $t:=\Abs{S}$ steps. Then, by using the vertices in $S_t\setminus S$, we can lower bound the number of edges in $H_{S_t}$:

    $$e(H_{S_t})\geq \Abs{S_t\setminus S} \e d= \frac{1}{2}\Abs{S_t}\e d\,.$$

    On the other hand, be \cref{lemma: emlonesided},

    $$e(H_{S_t})\leq \frac{\Abs{S_t}d}{2}\Paren{\frac{\Abs{S_t}}{n} + c}\,.$$

    Comparing the two and using $\Abs{S}=\frac{1}{2}\Abs{S_t}$ we get

    $$\Abs{S}\geq\frac{1}{2}\Paren{\e-c}n=\alpha n$$

    which contradicts upper bound assumption on $\Abs{S}$. Hence, letting $R:=V\setminus S_t$ it must hold that $\Abs{R}\geq n-2\Abs{S}\geq \Paren{1-2\alpha}n$ and the minimum degree in $H_R$ is at least $(1-\e)d=(1-2\alpha-c)d$ by construction.
\end{proof}


Now we are ready to prove \cref{lemma: approximatetopartial}.


\begin{proof}[Proof of \cref{lemma: approximatetopartial}]
    Let $c:=\frac{1}{9k}$. When $d$ is at least a large enough constant, $2^{-\Omega(d)}\leq \alpha$ holds. Let $B:=\Set{x\in V\mid \chi(x)\neq \hat{\chi}(x)}$ denote the set of vertices wrongly colored by $\hat{\chi}$. Note that $\Abs{B}\leq \alpha n$ by assumption, and $\Abs{SB(G)}\leq 2^{-\Omega(d)}\leq \alpha n$ by \cref{lemma: sbupperbound}. Hence, applying \cref{lemma: highmindegreesubgraph} on $H$ to $B\cup SB$ we get a set $R$ with $|R|=(1-4\alpha)n$ such that the minimum degree in $H_R$ is at least $\Paren{1-4\alpha-c}d$. Using $R\cap SB=\emptyset$, for any $x\in R$ and $c\neq \chi(x)$ we get that $x$ has degree at least $\Paren{\frac{1}{2k}-4\alpha-c}d\geq \frac{1}{6k}d$ into $\Set{y\in R\mid \chi(y)=c}$, i.e. into the vertices in $R$ colored with $c$. It follows by induction that no vertices in $R$ will be uncolored at any step of \cref{alg: uncoloring}. Hence, the coloring returned has size at least $|R|=(1-4\alpha)n$. It only remains to show that the coloring returned is proper. In fact, we will show that all vertices in $B$ will get uncolored, from which the statement follows.

    For the sake of contradiction, assume there exist vertices $B_2\subseteq B$ that remain colored by the end of the process. Let $x\in B_2$ be any such vertex. As $x$ has never been uncolored, at the end of the process it has at least $\frac{1}{6k}d$ neighbors $y$ with $\hat{\chi}(y)=c\neq \hat{\chi}(x)$. In particular, choosing $c=\chi(x)$, let $Q$ denote those neighbors $y$ of $x$ with $\hat{\chi}(y)=\chi(x)$. As $x$ and $y$ are neighbors in $G$, we have $\chi(x)\neq \chi(y)$, so $\hat{\chi}(y)\neq \chi(y)$, which implies $Q\subseteq B_2$. Summarizing, any vertex in $G_{B_2}$ has degree at least $\frac{1}{3k}d$. Therefore, $$e(G_{B_2})\geq \frac{\frac{1}{6k}d|B_2|}{2}= \frac{d|B_2|}{12k}\,.$$

    On the other hand, by \cref{lemma: emlonesided}, $$e(G_{B_2})\leq e(H_{B_2})\leq \frac{d\Abs{B_2}}{2}\Paren{\frac{\Abs{B_2}}{n}+c}\,.$$
    
    Using $\Abs{B_2}>0$ and the bounds on $\alpha$ and $c$, we can combine and simplify the above two inequalities to get $$\Abs{B_2}\geq \Paren{\frac{1}{6k}-c}n> \alpha n\,.$$

    However, as $\Abs{B_2}\leq \Abs{B}\leq \alpha n$ by assumption, we get a contradiction, so $\Abs{B_2}=0$ finishing the proof.
\end{proof}

Next we need an algorithm that recolors some of the uncolored vertices in a way that guarantees that the remaining uncolored vertices split up into small components.

\begin{algorithm}[H]
    As long as there is an uncolored vertex $x$ that has neighbors in every color except for some color $c$, color $x$ using $c$. 
    \protect\caption{\label{alg: saferecoloring} Safe Recoloring Algorithm}
\end{algorithm}

\begin{lemma}[Adaptation of Lemma B.21 of \cite{MR3536556-David16}]\label[lemma]{lemma: partialtologncomps}
    Let $H$ be a $d$-regular $n$-vertex  graph with $\lambda_2(H)\leq \frac{1}{16k^2}$. 
    Let $G$ be the graph obtained from $H$ by randomly planting a \kcl $\chi$ in it. Starting from a $\Paren{k, \frac{1}{256k^4}}$-partial coloring $\hat{\chi}$ of $G$, let $U$ be the set of vertices left uncolored by \cref{alg: saferecoloring}. Then, \whp each component in $G_U$ has size at most $\ln{n}$.
\end{lemma}

\begin{proof}
    The proof closely follows that of Lemma B.21 of \cite{MR3536556-David16}, so we only provide the parts that are changed. Let us first prove a helper lemma similar to Claim B.22 of \cite{MR3536556-David16}:
    \begin{lemma}\label[lemma]{lemma: largecoloredneighborhoodprob}
        Let $C$ be any set of vertices of size $t_1$, let $D:=N_{H}\Paren{C}\setminus C$
        and let $t_2:=\Abs{D}$. If $t_2\ge 2k^2t_1$ then $C\subseteq U$ and $D\subseteq V\setminus U$ with probability at most $e^{-\frac{1}{2k}t_2}$.
    \end{lemma}
    \begin{proof}
        Partition the vertices of $D$ to disjoint subsets $\Set{D_{v}}$, where $v\in C$ and $D_{v}\subseteq N_{H}\Paren{v}$. Consider some $D_{v}$ and assume that $v$ is colored by $1$. If both $v\in U$ and $D_{v}\subseteq V\setminus U$ then the set of colors that vertices in $D_{v}$ are using must be contained in exactly one
        of the sets $[k]\setminus {c}$ for some $c\neq 1$ (otherwise \cref{alg: saferecoloring} would color $v$).
        Therefore the probability that both $v\in U$ and $D_{v}\subseteq V\setminus U$
        is at most $(k-1)\Paren{\frac{k-1}{k}}^{\Abs{D_{v}}}$. Hence, the probability that both $C$ is in $U$ and $D\subseteq V\setminus U$ is at most
        \begin{align*}
            \prod_{D_{v}\in\Set{ D_{v}} }(k-1)\Paren{\frac{k-1}{k}}^{\Abs{D_{v}}} & =(k-1)^{\Abs{\Set{ D_{v}} }}\Paren{\frac{k-1}{k}}^{t_2}\\
            & \le(k-1)^{t_1}\Paren{\frac{k-1}{k}}^{t_2}\\
            %& =(k-1)^{t_1+t_2}k^{-t_2}\\
            & =e^{(t_1+t_2)\ln(k-1)-t_2\ln{k}}\\
            & \le e^{-\frac{1}{2k}t_2}
        \end{align*}
        using $t_2\ge 2k^2t_1$ and $\ln{k}-\ln(k-1)\geq \frac1{k}$ on the last line.
    \end{proof}

    Letting $N_{H}\Paren{n,t_1,t_2}$ denote the number of connected components of size $t_1$ with a neighborhood of size exactly $t_2$ and letting $P_{t_1,t_2}:=\max_{C\subset H,\Abs{C}=t_1,\Abs{N\Paren{C}}=t_2}\Pr\Brac{\text{\ensuremath{C}\,\ is in\,\ensuremath{U}}}$, we get by union bound that the probability that there exists a connected component in $G_{U}$ of size $\Omega(\ln n)$ is at most
    
    \begin{align}
    \sum_{t_1=\ln n\,}^{\Abs{U}}\sum_{t_2=1}^{d t_1}N_{H}\Paren{n,t_1,t_2}P_{t_1,t_2} & =\sum_{t_1=\ln n\,}^{\Abs{U}}\sum_{t_2= 7k^2t_1 }^{d t_1}N_{G}\Paren{n,t_1,t_2}e^{-\frac{1}{2k} t_2}\label{eq:alpha1}\\
    & \le\sum_{t_1=\ln n\,}^{\Abs{U}}\sum_{t_2= 7k^2t_1 }^{d t_1}n\binom{t_1+t_2 }{ t_1}e^{-\frac{1}{2k} t_2}\label{eq:alpha2}\\
    & \le\sum_{t_1=\ln n\,}^{\Abs{U}}\sum_{t_2= 7k^2t_1 }^{d t_1}n\Paren{e\frac{t_1+t_2}{t_1}}^{t_1}e^{-\frac{1}{2k} t_2}\label{eq:alpha3}\\
    & \le\sum_{t_1=\ln n\,}^{\Abs{U}}ndt_1\max_{7k^2\le x\le d}e^{\Paren{1+\ln(x+1)-\frac{1}{2k} x}t_1}\nonumber \\
    & \le n^{4}\max_{7k^2\le x\le d}e^{\Paren{1+\ln(x+1)-\frac{1}{2k} x}\ln{n}}\nonumber\\
    & \le n^{-1}\,.\label{eq:alpha4}
    \end{align}

    Equality~\ref{eq:alpha1} follows by \cref{lemma: largecoloredneighborhoodprob} and that by \cref{lemma: ssveonesided}
    if $\lambda_2(H)\le \frac{1}{16k^2}$ then for a set of size $t_1\le\Abs{U}\le\frac{1}{256k^4}n$ it holds that its neighborhood is of size $t_2\ge7k^2t_1$. Inequality~\ref{eq:alpha2} follows by Lemma B.23 of \cite{MR3536556-David16}. Inequality~\ref{eq:alpha3} follows by the binomial identity ${n \choose k}\le\Paren{e\frac{n}{k}}^{k}$. Inequality~\ref{eq:alpha4} follows by $\ln\Paren{7k^2+1}\leq \frac{7k}{2}-6$ which holds for $k\geq 3$. This completes the proof.
\end{proof}

We are now finally ready to prove \cref{thm: random-planting-full}:

\begin{proof}[Proof of \cref{thm: random-planting-full}]
    We can apply \cref{thm: random-planting-partial-list} to get a list of $O_{k}(1)$ $k$-colorings $\hat{\chi}: [n] \to [k]$ such that, for at least one of them, it holds up to a permutation of the color classes that $\mathbb{E}_{\bm{x} \sim \operatorname{Unif}([n])} \1[\hat\chi(\bm{x}) \neq \chi(\bm{x})] \leq O_k\Paren{1/d^{\Omega(1)}}$.

    Then, for $d$ at least a large enough constant compared to $k$, running \cref{alg: uncoloring} on all colorings in $L$, by \cref{lemma: approximatetopartial} it turns the coloring $\hat{\chi}$ with $\mathbb{E}_{\bm{x} \sim \operatorname{Unif}([n])} \1[\hat\chi(\bm{x}) \neq \chi(\bm{x})] \leq O_k\Paren{1/d^{\Omega(1)}}$ into a $\Paren{k, O_k\Paren{1/d^{\Omega(1)}}}$-partial coloring and returns a list of colorings $L'$. %\andor{should we remark that we somehow managed to avoid doing iterative recoloring which works only for two-sided expansion?}  

    Then, for $d$ at least a large enough constant compared to $k$, running \cref{alg: saferecoloring} on all colorings in $L'$, \cref{lemma: partialtologncomps} guarantees that, for the $\Paren{k, O_k\Paren{1/d^{\Omega(1)}}}$-partial coloring, all the remaining still uncolored vertices are in connected components of size $O(\ln{n})$ and returns a list of colorings $L''$:

    Hence, we can try a brute force search on each component of those colorings in $L''$ that have component sizes $O(\ln{n})$. By the above, we are guaranteed to find a full coloring on at least one of them. This finishes the proof.
\end{proof}
